%% parts/p08/paraconsistent-and-nonmonotonic.tex --- Part 8
\chapter{Paraconsistent and Nonmonotonic Logics}
\label{sec:p08-paraconsistent-and-nonmonotonic}

Paraconsistency changes what follows from inconsistent information;
nonmonotonicity changes how conclusions survive added information. These
properties are logically independent and must not be grouped merely because
both depart from classical consequence.

\section{Paraconsistent consequence and dialetheism}

\begin{hbdefinition}{def:p08-paraconsistency}{Paraconsistent consequence}
A consequence relation \(\vdash_L\) is paraconsistent when it is not explosive:
for some \(A,B\), \(A,\neg A\nvdash_L B\). The name is necessary to distinguish
failure of explosion from the stronger claim that some contradictions are
true. It stands as a property of a consequence relation, differs from
dialetheism's semantic/metaphysical thesis, and is used to classify LP,
relevance systems, and inconsistency-tolerant databases
\cite{Priest1979,Priest2006,OmoriWansing2017}.
\index{paraconsistent logic}\index{explosion}\index{dialetheism}
\end{hbdefinition}

Priest's LP uses the strong-Kleene truth functions with both true and both
true-and-false designated. Thus \(A\) and \(\neg A\) may both be designated
without every \(B\) being designated. K3 uses the same internal operations but
only true designated, so its consequence relation is different. Relevance
logics may block explosion for proof-theoretic or semantic reasons unrelated
to LP's designated set.

Dialetheism asserts that at least one contradiction is true. It entails a need
for nonexplosive consequence if reasoning is to remain nontrivial, but a
paraconsistent logician may deny every true contradiction and use the logic
only to reason safely about inconsistent information. Conversely, a many-valued
semantics is not paraconsistent unless its consequence relation actually fails
explosion.

\begin{hbdefinition}{def:p08-inclosure-schema}{Inclosure schema}
An \emph{inclosure instance} consists of a collection \(\Omega\), a condition
\(\delta(X)\) on subcollections \(X\subseteq\Omega\), and an operation
\(\psi\) such that
\[
 \delta(\Omega),\qquad
 \delta(X)\Longrightarrow\psi(X)\in\Omega,
 \qquad
 \delta(X)\Longrightarrow\psi(X)\notin X.
\]
Taking \(X=\Omega\) yields both
\(\psi(\Omega)\in\Omega\) and
\(\psi(\Omega)\notin\Omega\).  The name is necessary because the common
closure/transcendence pattern, rather than any one liar- or set-theoretic
sentence, is the object being classified.  It stands as a metalinguistic
paradox schema.  It differs from paraconsistent consequence, which controls
what follows after a contradiction is admitted, and is used to organize the
response families below \cite{Priest2006}.
\index{inclosure schema}\index{dialetheism!inclosure schema}
\end{hbdefinition}

Hierarchical and type-restrictive replies reject the closure step for the
unrestricted \(\Omega\); contextual and revisionary replies dispute the fixed
application of \(\delta\) or \(\psi\); consistency-preserving diagnoses reject
one of the two premises; dialetheic replies accept the resulting contradiction
but use a paraconsistent consequence relation to block explosion.  These are
different diagnoses of the displayed schema, not consequences of LP alone.

\begin{hbexample}{def:p08-lp-nonexplosion}{LP counterexample to explosion}
Assign \(A\) the value both and \(B\) false. Then \(A\) and \(\neg A\) are
designated while \(B\) is not, so \(A,\neg A\nvdash_{\mathsf{LP}}B\). This
single valuation proves nonexplosion; it does not establish that every LP
theory is nontrivial.
\end{hbexample}

\begin{hbdefinition}{def:p08-da-costa-c1}{da Costa's propositional calculus \(C_1\)}
To fix one member of da Costa's hierarchy rather than an undefined umbrella
``C-system,'' let \(\mathsf C_1\) be the exact propositional calculus on
\hyperref[sheet:hilbert-da-costa-c1]{\texttt{sheet:hilbert-da-costa-c1}}.
Its language has primitive \(\neg,\land,\lor,\to\), and
\[
 A^\circ\defeq\neg(A\land\neg A).
\]
The formula \(A^\circ\) is a formula-level consistency or well-behavedness
abbreviation in this calculus, not a matrix truth value and not the
metalinguistic assertion that the whole theory is consistent.  Its source
postulates restrict reductio by a \(B^\circ\) premise and propagate
well-behavedness through \(\land,\lor,\to\).

The system name and abbreviation are necessary because the da Costa hierarchy
contains distinct \(C_n\) calculi.  They stand here for this one Hilbert
calculus and formula abbreviation, differ from LP's designated-value
semantics, and are used in the nonexplosion theorem and comparison table below
\cite{daCosta1974}.
\index{da Costa calculus C1@da Costa calculus \(C_1\)}
\index{consistency formula!da Costa \(C_1\)}
\end{hbdefinition}

\begin{hbtheorem}{thm:p08-da-costa-c1-nonexplosion}{cited=daCosta1974}{Nonexplosion and controlled explosion in \(C_1\)}
Unrestricted explosion fails:
\[
 A,\neg A\nvdash_{\mathsf C_1}B
\]
for suitable formulas \(A,B\); equivalently,
\((A\land\neg A)\to B\) is not a theorem schema.  But the consistency formula
restores controlled explosion:
\[
 A^\circ,A,\neg A\derives_{\mathsf C_1}B.
\]
The second judgment does not contradict the first because \(A^\circ\) is an
additional object-language premise.
\end{hbtheorem}

\begin{table}[htbp]
\centering\small
\begin{tabularx}{\textwidth}{@{}l X X X@{}}
\toprule
System & Exact formal mechanism & Explosion boundary & Standing \\
\midrule
\(\mathsf{LP}\) & strong-Kleene three-value matrix with both and true
designated & \(A,\neg A\nvdash_{\mathsf{LP}}B\) by
Example~\ref{def:p08-lp-nonexplosion}; no consistency connective belongs to
the fixed LP language & semantic consequence by designated-value preservation
\cite{Priest1979} \\
\(\mathsf C_1\) & positive-classical Hilbert core with paraconsistent
negation and \(A^\circ=\neg(A\land\neg A)\) & unrestricted explosion fails,
but \(A^\circ,A,\neg A\derives_{\mathsf C_1}B\) & proof-theoretic
consistency control on
\hyperref[sheet:hilbert-da-costa-c1]{the exact \(C_1\) sheet}
\cite{daCosta1974} \\
\bottomrule
\end{tabularx}
\caption{LP and da Costa \(C_1\): two nonexplosive mechanisms.  LP controls
consequence through designation; \(C_1\) internalizes a formula-level
consistency condition that licenses classical recovery locally.}
\label{tab:p08-paraconsistent-systems}
\end{table}

\section{Nonmonotonic consequence}

\begin{hbdefinition}{def:p08-nonmonotonic}{Nonmonotonic consequence}
A consequence relation \(\mathrel{|\!\sim}\) is nonmonotonic when it may have
\(\Gamma\mathrel{|\!\sim}A\) but
\(\Gamma\cup\{B\}\not\mathrel{|\!\sim}A\). The notation is needed to
distinguish defeasible inference from the monotonic semantic and proof
relations used elsewhere. It stands as a relation between information states
and conclusions, differs because additional premises may defeat a conclusion,
and is used in preferential, default, and belief-revision systems
\cite{KrausLehmannMagidor1990,Makinson2005}.
\index{nonmonotonic consequence}\index{defeasible inference}
\end{hbdefinition}

\begin{hbdefinition}{def:p08-preferential-model}{Smooth preferential and ranked models}
A \emph{preferential model} is a triple \((S,\ell,\prec)\), where \(S\) is a
set of states, \(\ell(s)\) is a classical propositional valuation, and
\(\prec\) is a strict partial order: it is irreflexive and transitive.  For
\(X\subseteq S\), put
\[
 \mu(X)=\{s\in X:\text{there is no }t\in X\text{ with }t\prec s\},
 \qquad
 [A]=\{s:\ell(s)\models A\}.
\]
The model is \emph{smooth} when every \(s\in X\) is either in \(\mu(X)\) or
lies above some member of \(\mu(X)\), for every formula-definable \(X\).
It represents defeasible consequence by
\(A\mathrel{|\!\sim}B\) exactly when
\(\mu([A])\subseteq[B]\).  It is \emph{ranked} when a map into a totally
ordered set of layers satisfies
\(s\prec t\) exactly when the rank of \(s\) is lower than the rank of \(t\).

These names are necessary because mere existence of minimal worlds does not
give the representation theorem.  They name semantic structures, not new
inference rules.  Smoothness differs from ranking: the former grounds System
P, while the latter imposes the stronger layered structure used for Rational
Monotony.  Both distinctions are used immediately below
\cite{KrausLehmannMagidor1990,LehmannMagidor1992}.
\index{preferential model}\index{smoothness!preferential model}
\index{ranked model}
\end{hbdefinition}

\begin{table}[htbp]
\centering\small
\begin{tabularx}{\textwidth}{@{}l X X@{}}
\toprule
System-P principle & Premises & Conclusion \\
\midrule
Supraclassicality & \(A\vdash_{\mathsf{CPL}}B\) &
  \(A\mathrel{|\!\sim}B\) \\
Left logical equivalence & \(\vdash A\leftrightarrow B\),
  \(A\mathrel{|\!\sim}C\) & \(B\mathrel{|\!\sim}C\) \\
Right weakening & \(\vdash B\to C\), \(A\mathrel{|\!\sim}B\) &
  \(A\mathrel{|\!\sim}C\) \\
And & \(A\mathrel{|\!\sim}B\), \(A\mathrel{|\!\sim}C\) &
  \(A\mathrel{|\!\sim}B\land C\) \\
Or & \(A\mathrel{|\!\sim}C\), \(B\mathrel{|\!\sim}C\) &
  \(A\lor B\mathrel{|\!\sim}C\) \\
Cut & \(A\mathrel{|\!\sim}B\), \(A\land B\mathrel{|\!\sim}C\) &
  \(A\mathrel{|\!\sim}C\) \\
Cautious Monotony & \(A\mathrel{|\!\sim}B\),
  \(A\mathrel{|\!\sim}C\) & \(A\land B\mathrel{|\!\sim}C\) \\
Rational Monotony & \(A\mathrel{|\!\sim}C\),
  \(A\not\mathrel{|\!\sim}\neg B\) &
  \(A\land B\mathrel{|\!\sim}C\) \\
\bottomrule
\end{tabularx}
\caption{System P and its rational extension.  Rational Monotony is not a
System-P rule; it is displayed to mark the ranked-model extension.}
\label{tab:p08-system-p-rules}
\end{table}

\begin{hbtheorem}{thm:p08-klm-representation}{cited=KrausLehmannMagidor1990}{KLM representation for System P}
For the propositional language and formula-definable sets fixed above, a
defeasible consequence relation satisfies the seven System-P principles in
Table~\ref{tab:p08-system-p-rules} exactly when it is represented by a smooth
preferential model.  The smoothness hypothesis is part of the theorem.
\end{hbtheorem}

\begin{hbtheorem}{thm:p08-rational-representation}{cited=LehmannMagidor1992}{Ranked representation for rational consequence}
A System-P relation satisfying Rational Monotony has a ranked-model
representation, and every ranked model validates Rational Monotony.  This
stronger conclusion is not inferred from an arbitrary smooth preference.
\end{hbtheorem}

\begin{table}[htbp]
\centering
\scriptsize
\begin{tabularx}{\textwidth}{@{}l X X X@{}}
\toprule
Framework & Semantic/operational object & Consequence selected & Finite propositional complexity \\
\midrule
preferential consequence & smooth model
\((S,\ell,\prec)\) & truth at every minimal premise-state & model-relative evaluation; representation is Theorem~\ref{thm:p08-klm-representation} \\
ranked consequence & totally ordered exception layers & truth at every lowest-rank premise-state & model-relative evaluation; rational representation is Theorem~\ref{thm:p08-rational-representation} \\
Reiter default logic & grounded fixed points \(E=\Gamma_E\) & credulous: some extension; skeptical: all extensions & extension existence and credulous reasoning \(\Sigma_2^p\)-complete; skeptical reasoning \(\Pi_2^p\)-complete \\
circumscription & models minimal in the declared \(\bar P\), with fixed/varying symbols typed & truth in every minimal model & propositional inference \(\Pi_2^p\)-complete \\
autoepistemic logic & stable expansions of Definition~\ref{def:p08-autoepistemic-expansion} & credulous: some expansion; skeptical: all expansions & expansion existence/credulous reasoning \(\Sigma_2^p\)-complete; skeptical reasoning \(\Pi_2^p\)-complete \\
AGM revision & operation \(K*A\) on a deductively closed belief set & membership in the one revised theory & depends on the finite representation of \(K\) and \(*\); no representation-free class asserted \\
\bottomrule
\end{tabularx}
\caption{Nonmonotonic frameworks, consequence selection, and standard finite
propositional decision boundaries.  Default, circumscription, and
autoepistemic bounds follow \textcite{Gottlob1992}.}
\label{tab:p08-nonmonotonic-properties}
\end{table}

\section{Belief revision}

\begin{hbdefinition}{def:p08-agm}{AGM revision and contraction}
Let \(K\) be a deductively closed classical belief set. A revision operation
\(K*A\) maps a formula input to a deductively closed belief set.  Write
\(K+A=\operatorname{Th}(K\cup\{A\})\) for expansion.  A contraction operation
\(K-A\) maps \(A\) to a deductively closed subset of \(K\) intended no longer
to imply \(A\), except when \(A\) is a tautology.  The revision and contraction
postulates are displayed in
Tables~\ref{tab:p08-agm-postulates}
and~\ref{tab:p08-agm-contraction-postulates}.  In particular, revision
Inclusion is unconditional, whereas the consistency hypothesis belongs to
Vacuity.  These names are necessary to distinguish adding an input
consistently, removing commitment to it, and merely expanding by it.  They
stand as operations on belief sets, differ from nonmonotonic consequence by
taking an explicit prior state and input, and are used by the identities and
representation theorem below
\cite{AlchourronGardenforsMakinson1985}.
\index{AGM belief revision}\index{belief contraction}\index{belief update}
\end{hbdefinition}

\begin{table}[htbp]
\centering\small
\begin{tabularx}{\textwidth}{@{}l X@{}}
\toprule
Postulate & Exact revision condition \\
\midrule
Closure & \(K*A=\operatorname{Th}(K*A)\) \\
Success & \(A\in K*A\) \\
Inclusion & \(K*A\subseteq K+A\) \\
Vacuity & if \(\neg A\notin K\), then \(K+A\subseteq K*A\) \\
Consistency & if \(A\) is classically consistent, then \(K*A\) is consistent \\
Extensionality & if \(\vdash A\leftrightarrow B\), then \(K*A=K*B\) \\
Superexpansion & \(K*(A\land B)\subseteq (K*A)+B\) \\
Subexpansion & if \(\neg B\notin K*A\), then
  \((K*A)+B\subseteq K*(A\land B)\) \\
\bottomrule
\end{tabularx}
\caption{The eight AGM postulates for one-step revision.}
\label{tab:p08-agm-postulates}
\end{table}

\begin{table}[htbp]
\centering\small
\begin{tabularx}{\textwidth}{@{}l X@{}}
\toprule
Postulate & Exact contraction condition \\
\midrule
Closure & \(K-A=\operatorname{Th}(K-A)\) \\
Inclusion & \(K-A\subseteq K\) \\
Vacuity & if \(A\notin K\), then \(K-A=K\) \\
Success & if \(\not\vdash A\), then \(A\notin K-A\) \\
Recovery & \(K\subseteq (K-A)+A\) \\
Extensionality & if \(\vdash A\leftrightarrow B\), then \(K-A=K-B\) \\
Conjunctive overlap &
  \((K-A)\cap(K-B)\subseteq K-(A\land B)\) \\
Conjunctive inclusion & if \(A\notin K-(A\land B)\), then
  \(K-(A\land B)\subseteq K-A\) \\
\bottomrule
\end{tabularx}
\caption{The eight AGM postulates for belief-set contraction.  Success excludes
tautological inputs; the final two rows are the supplementary conjunctive
postulates.}
\label{tab:p08-agm-contraction-postulates}
\end{table}

The two cross-operation identities are
\[
 K*A=(K-\neg A)+A
 \quad\text{(Levi)},
 \qquad
 K-A=K\cap(K*\neg A)
 \quad\text{(Harper)}.
\]
Partial-meet constructions choose maximal subsets of \(K\) that do not imply
\(A\), then intersect selected remainders.  Epistemic-entrenchment and
faithful preorder representations characterize equivalent presentations under
their respective representation hypotheses
\cite{AlchourronGardenforsMakinson1985}.

\begin{hbdefinition}{def:p08-agm-faithful-preorder}{Faithful preorder for a belief set}
Fix a finite propositional signature and let \(\mathcal W\) be its valuations.
For a consistent belief set \(K\), a total preorder \(\preceq_K\) on
\(\mathcal W\) is \emph{faithful to \(K\)} exactly when all valuations
satisfying \(K\) are tied and every such valuation is strictly preferred to
every valuation that does not satisfy \(K\).  Write
\(\min([A],\preceq_K)\) for the \(\preceq_K\)-minimal \(A\)-valuations and
\(\operatorname{Th}(X)\) for the formulas true at every member of \(X\).

The qualifier is necessary because an arbitrary total preorder need not
preserve the prior beliefs when they are consistent with the input.  It names
semantic data assigned to \(K\), not another revision operation, and differs
from the preferential models of Definition~\ref{def:p08-preferential-model}
because it orders all classical valuations relative to one fixed prior belief
set.  The representation theorem uses it immediately
\cite{KatsunoMendelzon1991}.
\end{hbdefinition}

\begin{hbtheorem}{thm:p08-agm-representation}{cited=KatsunoMendelzon1991}{AGM faithful-preorder representation for finite propositional belief sets}
For each consistent \(K\), a revision operation satisfies all eight postulates
in Table~\ref{tab:p08-agm-postulates} exactly when there is a total preorder
\(\preceq_K\) faithful to \(K\) such that, for every formula \(A\),
\[
 K*A=\operatorname{Th}\bigl(\min([A],\preceq_K)\bigr).
\]
If \(A\) is inconsistent, the minimum set is empty and its theory is the
inconsistent theory, as required by the consistency postulate.  The original
AGM partial-meet representation instead constructs contraction and obtains
revision through Levi; the two presentations agree under their stated
propositional representation hypotheses
\cite{AlchourronGardenforsMakinson1985}.  Database update, iterated revision,
and belief-base revision require different or additional postulates and are not
included.
\end{hbtheorem}

\section{Default logic and circumscription}

\begin{hbdefinition}{def:p08-default-theory}{Reiter default theory}
A default is written
\(\frac{\alpha:\beta_1,\ldots,\beta_n}{\gamma}\): prerequisite \(\alpha\),
justifications \(\beta_i\), and conclusion \(\gamma\). A default theory
\((W,D)\) has facts \(W\) and defaults \(D\).  For a candidate theory \(E\),
let \(\Gamma_E\) be the least set \(U\) satisfying:
\begin{enumerate}
  \item \(W\subseteq U\);
  \item \(U=\operatorname{Th}(U)\); and
  \item if
    \(\frac{\alpha:\beta_1,\ldots,\beta_n}{\gamma}\in D\),
    \(\alpha\in U\), and \(\neg\beta_i\notin E\) for every \(i\), then
    \(\gamma\in U\).
\end{enumerate}
Equivalently, start with \(U_0=\operatorname{Th}(W)\), close each successor
stage under classical consequence after adding every conclusion whose
prerequisite is in the current stage and whose justifications are consistent
with \(E\), and take the union of the stages.  The candidate \(E\) is a
\emph{Reiter extension} exactly when \(E=\Gamma_E\).

The term is needed because applicability tests justifications against the
candidate \(E\), while prerequisite support is grounded in the least set
\(\Gamma_E\).  It stands as a theory-plus-fixed-point semantics, differs from
the circular condition that tests prerequisites directly in \(E\), and is used
to explain absent, multiple, credulous, and skeptical conclusions
\cite{Reiter1980}.
\index{default logic}\index{default extension}
\end{hbdefinition}

A default theory may have no extension or several. Credulous consequence holds
in some extension; skeptical consequence holds in all. The normal default
\(\frac{Bird(x):Fly(x)}{Fly(x)}\) supports flight until information
\(\neg Fly(x)\) defeats its justification. This is not material implication.

\begin{hbdefinition}{def:p08-circumscription}{Predicate circumscription}
Write the first-order theory as \(T(\bar P,\bar Q,\bar Z)\), where \(\bar P\)
is minimized, \(\bar Q\) is fixed, and \(\bar Z\) may vary.  For predicate
tuples of matching arities, let \(\bar P'\leq\bar P\) abbreviate the
componentwise universal closures of \(P'_i(\bar x)\to P_i(\bar x)\).
Circumscription is the second-order sentence
\[
 T(\bar P,\bar Q,\bar Z)\land
 \forall\bar P'\forall\bar Z'
 \Bigl((T(\bar P',\bar Q,\bar Z')\land\bar P'\leq\bar P)
       \to\bar P\leq\bar P'\Bigr).
\]
Thus no model with the same fixed \(\bar Q\) can make the minimized predicates
strictly smaller while adjusting only the declared varying predicates.  The
name is necessary because the minimized, fixed, and varying tuples determine
different policies.  It stands as a model-minimization transformation,
differs from syntactic default application, and is used to formalize minimal
abnormality \cite{McCarthy1980}.
\index{circumscription}\index{predicate minimization}
\end{hbdefinition}

Parallel circumscription minimizes several predicates together; prioritized
circumscription minimizes them in an order. Allowing a predicate to vary can
change the minimal models even when it is not minimized. These parameters must
be recorded before a circumscription consequence is compared with default or
preferential consequence.

\begin{hbdefinition}{def:p08-autoepistemic-expansion}{Stable expansion in autoepistemic logic}
Let \(T\) be a theory in the Boolean language closed under the belief operator
\(L\).  Write \(\operatorname{Th}\) for ordinary classical propositional
closure over that language: modal-prefixed subformulas are formulas, but this
closure adds no modal axiom or necessitation rule.  A \emph{stable expansion}
is a theory \(E\) satisfying
\[
 E=\operatorname{Th}
 \bigl(T\cup\{LA:A\in E\}\cup\{\neg LA:A\notin E\}\bigr).
\]
The name is necessary because the agent's positive and negative introspective
assumptions are determined by the candidate theory itself.  It names a fixed
point, not a Reiter default extension: default prerequisites are grounded
through \(\Gamma_E\), whereas the belief and nonbelief assumptions here are
closed as one expansion.  Its classical closure also differs from simply
declaring \(L\) to be a KD45 modality; the comparison table uses this exact
distinction \cite{Moore1985}.
\index{autoepistemic logic}\index{stable expansion!autoepistemic}
\end{hbdefinition}

\begin{systementry}{entry:p08-reiter-default}{Reiter propositional default logic}
  \entrylanguage{Classical propositional formulas plus defaults
    \(\frac{\alpha:\beta_1,\ldots,\beta_n}{\gamma}\).}
  \entrysemantics{Grounded extensions \(E=\Gamma_E\),
    Definition~\ref{def:p08-default-theory}.}
  \entrycalculus{sheet:reiter-default-extension}
  \entryadequacy{cited=Reiter1980}{The operational sheet is the displayed
    least grounded extension construction; it does not admit circular
    self-support.}
  \entrydecidability{cited=Gottlob1992}{Finite propositional extension
    existence and credulous inference are \(\Sigma_2^p\)-complete; skeptical
    inference is \(\Pi_2^p\)-complete.}
  \entrycharacteristic{Justifications are checked against the candidate
    extension, while prerequisites must arise in the grounded stages.}
  \entrytranslations{No material-implication translation preserves the
    extension operator without additional bookkeeping.}
  \entrysources{Reiter1980,Gottlob1992}
\end{systementry}

\begin{systementry}{entry:p08-circumscription}{Predicate circumscription}
  \entrylanguage{Second-order classical logic over minimized \(\bar P\), fixed
    \(\bar Q\), and varying \(\bar Z\) predicate tuples.}
  \entrysemantics{Models minimal in \(\bar P\) under the exact second-order
    sentence of Definition~\ref{def:p08-circumscription}.}
  \entrycalculus{sheet:sol}
  \entryadequacy{cited=McCarthy1980}{Second-order consequence from the displayed
    circumscription sentence agrees with minimal-model consequence for the
    declared fixed/varying policy.}
  \entrydecidability{cited=Gottlob1992}{Propositional circumscription inference
    is \(\Pi_2^p\)-complete.}
  \entrycharacteristic{Changing a predicate from fixed to varying can change
    the minimal models even when it is not minimized.}
  \entrytranslations{Parallel and prioritized circumscription require distinct
    minimization policies.}
  \entrysources{McCarthy1980,Gottlob1992}
\end{systementry}

\begin{systementry}{entry:p08-autoepistemic}{Moore autoepistemic logic}
  \entrylanguage{The propositional modal language with belief operator \(L\).}
  \entrysemantics{Stable expansions of
    Definition~\ref{def:p08-autoepistemic-expansion}.}
  \entrycalculus{sheet:hilbert-cpl}
  \entryadequacy{cited=Moore1985}{Classical closure in the autoepistemic
    language plus the displayed stable-expansion operator fixes the source
    presentation; no modal axiom is silently added to \(\operatorname{Th}\).}
  \entrydecidability{cited=Gottlob1992}{Finite propositional expansion
    existence and credulous inference are \(\Sigma_2^p\)-complete; skeptical
    inference is \(\Pi_2^p\)-complete.}
  \entrycharacteristic{Positive and negative introspective assumptions are
    determined by membership and nonmembership in the candidate expansion.}
  \entrytranslations{Stable expansions and Reiter extensions are neighboring
    fixed-point formalisms, not identical objects.}
  \entrysources{Moore1985,Gottlob1992}
\end{systementry}

\begin{hbwarning}{def:p08-nonmonotonic-status}{Nonmonotonic results are framework-specific}
Failure of monotonicity alone proves no cumulativity, rationality, existence of
extensions, or computational bound. Every claim in this chapter fixes its
underlying preference, ranking, default, minimization, or revision object;
translations between frameworks require separate preservation
theorems.
\end{hbwarning}
